\begin{abcliste}
\abc
In a harmonic oscillation the acceleration is proportional to the displacement, $a = -\omega^2 y$. In the air, however, the ball accelerates with $g$ downwards at every height: $a = -g$, whatever the height. Its graph is a row of parabolas, not a cosine. The motion is \result{periodic, but not harmonic.}

Losing energy at each bounce is not the reason: even a ball that bounced without any losses would not oscillate harmonically.

\abc
Up and down are a free fall each: falling from the top takes $t = \sqrt{2h/g}$, rising takes as long. Hence
\begin{align}
 T &= \TF \\
   &= 2\sqrt{\frac{2\times\hB}{\ncg}} \\
   &= \T \approx \result{\TP{0}{2}}
\end{align}

\abc
\begin{align}
 T_4 &= \TbF = 2\,T \\
   &= \Tb \approx \result{\TbP{0}{2}}
\end{align}
Four times as high, twice as long: the period depends on the amplitude, unlike in a harmonic oscillation.
\end{abcliste}
